\section{利用微分三角形求曲线的弧长}
\markboth{12. 利用微分三角形求曲线的弧长}{12. 利用微分三角形求曲线的弧长}

在讨论曲线的性质和应用时，很多问题都与曲线的弧长有关，这时可以利用下面介绍的微分三角形。

如图 12.1，设 $M(x, y)$ 是曲线 $\Gamma$ 上的任意点，$N(x + dx, y + \Delta y)$ 是曲线 $\Gamma$ 上靠近 $M$ 的点。为了看得清楚，图中把 $N$ 与 $M$ 之间的距离有意画大了。又设图中 $NP$ 平行于 $y$ 轴，$MP$ 平行于 $x$ 轴，$MQ$ 是曲线 $\Gamma$ 在 $M$ 点的切线，$MQ$ 与 $x$ 轴的夹角为 $\alpha$，那么从 $MP = dx$ 和
\[
  \tg\alpha = \frac{dy}{dx}
\]
可得
\[
  PQ = dy,\quad PN = \Delta y = dy + QN.
\]

\begin{center}
  \begin{minipage}{\linewidth}
    \centering
    \begin{tikzpicture}[scale=0.90, >=Stealth, font=\footnotesize]
      % 坐标轴
      \draw[->] (-0.45, 0) -- (3.65, 0) node[below left] {$x$};
      \draw[->] (0, -0.3) -- (0, 3.15) node[below left] {$y$};
      \node[below left] at (0, 0) {$O$};

      % 曲线 Gamma: y = 0.5 + 0.30*x + 0.16*x^2
      \draw[very thick, curveblue, smooth, domain=-0.2:2.95, samples=80]
        plot (\x, {0.5 + 0.30*\x + 0.16*\x*\x});
      \node[curveblue, left] at (-0.08, 0.58) {$\Gamma$};

      % 关键点坐标
      \coordinate (M) at (0.65, 0.7626);
      \coordinate (N) at (2.65, 2.4186);
      \coordinate (P) at (2.65, 0.7626);
      \coordinate (Q) at (2.65, 1.7786);
      \coordinate (C) at (0.65, 0);
      \coordinate (D) at (2.65, 0);
      \coordinate (T) at (3.00, {0.7626 + 0.508*(3.00 - 0.65)});

      % 辅助垂线与水平线
      \draw[dashed, mutedgray] (M) -- (C) node[below, black] {$C$};
      \draw[dashed, mutedgray] (P) -- (D) node[below, black] {$D$};
      \draw[thin] (M) -- (P);
      \draw[thin] (P) -- (N);

      % 切线 MQ（略微延长）
      \draw[thick] (M) -- (T);

      % 直角标记
      \draw ($(P)+(-0.15,0)$) -- ($(P)+(-0.15,0.15)$) -- ($(P)+(0,0.15)$);

      % 夹角 alpha
      \pic[draw, angle radius=0.44cm, angle eccentricity=1.35, "\scriptsize$\alpha$"] {angle=P--M--Q};

      % 点标记
      \fill (M) circle (1.4pt) node[above left, inner sep=1.5pt] {$M$};
      \fill (N) circle (1.4pt) node[left, inner sep=2pt] {$N$};
      \fill (P) circle (1.4pt) node[right, inner sep=1.5pt] {$P$};
      \fill (Q) circle (1.4pt) node[below right, xshift=-1pt, yshift=1pt] {$Q$};

      % 边长与增量标注
      \node[below] at ($(M)!0.5!(P)$) {$dx$};
      \node[left] at ($(P)!0.5!(Q)$) {$dy$};
      \node[below right, inner sep=1pt] at ($(M)!0.5!(Q)$) {$ds$};
      \node[above left, inner sep=2pt] at (1.55, 1.48) {$\Delta s$};

      % 右侧大括号标注 Delta y
      \draw[decorate, decoration={brace, amplitude=4.5pt, raise=16pt}]
        (N) -- (P) node[midway, right=22pt] {$\Delta y$};
    \end{tikzpicture}
    \captionof*{figure}{图 12.1}
  \end{minipage}
\end{center}

当 $dx$ 非常小时，$QN$ 比 $dy$ 小得多，可以忽略，这时可以近似地认为 $N$ 点与 $Q$ 点重合。相应地，一小段曲线弧 $\arc{MN}$ 可以近似看成切线上的对应线段 $MQ$。用 $\Delta s$ 表示弧 $\arc{MN}$ 的长度，$ds$ 表示线段 $MQ$ 的长度，那么当 $dx$ 非常小时就有
\[
  \Delta y \approx dy,\quad \Delta s \approx ds.
\]
因而，正像 $dy$ 叫做 $y$ 的微分一样，$ds$ 叫做弧长 $s$ 的微分。对于直角三角形 $MPQ$，应用勾股定理得到
\[
  (ds)^2 = (dx)^2 + (dy)^2.
\]
通常把上式简写成
\[
  ds^2 = dx^2 + dy^2,
\]

\noindent
\begin{minipage}[c]{0.38\textwidth}
  \centering
  \begin{tikzpicture}[scale=0.76, font=\footnotesize]
    \coordinate (M) at (0, 0);
    \coordinate (P) at (3.0, 0);
    \coordinate (N) at (3.0, 1.95);

    \draw[thick] (M) -- (P) -- (N);
    \draw[very thick, curveblue] (M) -- (N);

    % 直角标记
    \draw ($(P)+(-0.18,0)$) -- ($(P)+(-0.18,0.18)$) -- ($(P)+(0,0.18)$);
    % 夹角 alpha
    \pic[draw, angle radius=0.45cm, angle eccentricity=1.35, "\scriptsize$\alpha$"] {angle=P--M--N};

    % 顶点与边标注
    \node[below left] at (M) {$M$};
    \node[below right] at (P) {$P$};
    \node[above right] at (N) {$N$};
    \node[below] at ($(M)!0.5!(P)$) {$dx$};
    \node[right] at ($(P)!0.5!(N)$) {$dy$};
    \node[above left] at ($(M)!0.5!(N)$) {$ds$};
  \end{tikzpicture}
  \captionof*{figure}{图 12.2}
\end{minipage}\hfill
\begin{minipage}[c]{0.58\textwidth}
  \noindent 并且把线段 $MQ$ 画成与曲线弧 $\arc{MN}$ 重合在一起，使图 12.1 简化成图 12.2 的形状。图 12.2 中由微分 $dx$、$dy$ 和 $ds$ 组成的直角三角形就叫做微分三角形。微分三角形很容易记住，并且从它立刻可写出
  \[
    ds = \sqrt{dx^2 + dy^2}.
  \]
\end{minipage}

$ds$ 只不过是一小段弧长。再将 $ds$ 积分，就能求出弧长 $s$ 了。具体说来，如果曲线是由参数方程 $x = \varphi(\theta)$，$y = f(\theta)$ 给出的，用撇号表示对参数 $\theta$ 的导数，那么
\[
  ds = \sqrt{(x'\dd\theta)^2 + (y'\dd\theta)^2} = \sqrt{x'^2 + y'^2}\dd\theta,
\]
因而曲线在 $\theta = \theta_1$ 和 $\theta = \theta_2$ 的两点间的弧长是
\[
  s = \int_{\theta_1}^{\theta_2} \sqrt{x'^2 + y'^2}\dd\theta.
\]

直角坐标方程 $y = f(x)$ 可看成参数方程的特例，这时参数 $\theta = x$，因而 $x' = 1$，曲线在 $x = x_1$ 和 $x = x_2$ 的两点间的弧长是
\[
  s = \int_{x_1}^{x_2} \sqrt{1 + y'^2}\dd x.
\]

作为例子，我们现在来求摆线一拱的长。摆线方程是
\[
  x = r(\theta - \sin\theta),\quad y = r(1 - \cos\theta).\quad (0 \leqslant \theta \leqslant 2\pi)
\]
由此得到摆线一拱的弧长是
\begin{align*}
  s &= \int_{0}^{2\pi} \sqrt{x'^2 + y'^2}\dd\theta \\
    &= \int_{0}^{2\pi} \sqrt{r^2(1 - \cos\theta)^2 + r^2\sin^2\theta}\dd\theta \\
    &= \int_{0}^{2\pi} r\sqrt{1 - 2\cos\theta + \cos^2\theta + \sin^2\theta}\dd\theta \\
    &= \int_{0}^{2\pi} r\sqrt{2(1 - \cos\theta)}\dd\theta
     = \int_{0}^{2\pi} r\sqrt{4\sin^2\frac{\theta}{2}}\dd\theta \\
    &= 2r\int_{0}^{2\pi} \sin\frac{\theta}{2}\dd\theta
     = \left.-4r\cos\frac{\theta}{2}\right|_{0}^{2\pi} = 8r.
\end{align*}
由于摆线的拱高是 $2r$，所以摆线一拱的弧长恰好是拱高的 $4$ 倍。
